\documentclass[12pt]{article}
\usepackage[utf8]{inputenc}
\usepackage[T1]{fontenc}
\usepackage{hyperref}
\title{Trumai Language Resources for K-12 Curriculum Development}
\author{Compiled by Claude Sonnet 5 (claude-sonnet-5) for Andr\'as Kornai}
\date{}
\begin{document}
\maketitle

\section{Introduction}
Trumai (ISO 639-3 tpy; Glottocode trum1247) is a language isolate spoken in the Xingu Indigenous Park, Brazil. UNESCO classifies it as critically endangered. Wikipedia gives 120 ethnic Trumai and 51 speakers, both from 2006. Most Trumai speak languages of wider communication, and children are not learning Trumai well. The only substantial resources found are academic.

\section{Primary Educational Resources}
I found no Trumai primers, textbooks or school programmes.

\section{Linguistic and Reference Materials}
\begin{itemize}
\item Raquel Guirardello (1999), \emph{A Reference Grammar of Trumai}, PhD thesis, Rice University. Listed at 475 pages. The Socioambiental record does not link a full text. \url{https://acervo.socioambiental.org/acervo/tesesdissertacoes/reference-grammar-trumai}
\item The same record on a second Socioambiental address: \url{https://acervo-tmp.socioambiental.org/acervo/tesesdissertacoes/reference-grammar-trumai}
\item Guirardello-Damian, ``Trumai'', a chapter with a myth in IPA transcription, glosses and a free translation (Langsci Press): \url{https://langsci-press.org/catalog/view/167/741/840-1}
\item ``Ergativity in Trumai'' (2010), in \emph{Ergativity in Amazonia}, ed. Gildea and Queixal\'os (Benjamins): \url{https://www.jbe-platform.com/content/books/9789027288509-tsl.89.09raq}
\item Chapter page for the same volume: \url{https://lwc1.benjamins.com/catalog/tsl.89.09raq}
\item Article in Revista Linguistica (UFRJ), which I did not open: \url{https://revistas.ufrj.br/index.php/rl/article/download/25562/13842/62713}
\item Wikipedia overview of the language: \url{https://en.wikipedia.org/wiki/Trumai_language}
\end{itemize}

\section{Religious and Community Resources}
The search found none.

\section{Digital and Audio Resources}
The Langsci Press text chapter above is the only text sample found. I found no audio or video resources.

\section{Curriculum Development Context}
With about 51 or fewer speakers, mostly older, Trumai needs documentation and revitalization work before any K-12 curriculum is possible. The Guirardello grammar and texts are the main starting point. Any work should be done with the Trumai community in the Xingu Park.
\end{document}